\studySubsection{calc-13-multivariable-differential-calculus}{第 13 讲 多元函数微分学}

\begin{knowledgeBox}
教材页码：第 304 页起。

本讲用于整理多元函数极限、连续、偏导、全微分、方向导数、梯度、隐函数和多元极值。新增题目时重点区分连续、偏导存在、可微之间的关系。
\end{knowledgeBox}

\studySubsection{calc-13-k132}{K132 多元函数概念、定义域与图形}
\knowledgeAnchor[MATH1-KN-CALC-0179]{多元函数概念、定义域与图形}

\begin{knowledgeBox}
\textbf{定位与前置桥。}
一元函数把一个自变量映到一个函数值；二元函数
\(z=f(x,y)\) 则把平面区域 \(D\subset\mathbb R^2\) 中的每个点
\((x,y)\) 唯一映到一个实数。这里“区域”是允许输入的点集，
“图形”是三维空间中的点集
\(\{(x,y,f(x,y)):(x,y)\in D\}\)。本节需要的一元前置只有：
根号内非负、对数真数为正、分母不为零。它们在多元情形中仍逐条成立，
但最后必须取所有约束的交集；例如
\(\sqrt{1-x^2-y^2}\) 的最小模型是闭圆盘，而不是两个独立区间。

\textbf{严格核心与流程。}
求定义域时依次列出每个表达式的合法条件、求交并标明边界是否包含；
解释图形时区分“定义域在 \(xy\) 平面”与“函数图形在 \(xyz\) 空间”。
耦合条件如 \(x^2+y^2<1\) 不能拆成 \(|x|<1,|y|<1\)，因为后者是正方形。

\textbf{例1（闭边界）。}
\emph{信号：}根式 \(f=\sqrt{4-x^2-y^2}\)。
\emph{分析与解答：}须有 \(x^2+y^2\le4\)，故定义域为含边界的半径 \(2\)
闭圆盘，图形满足 \(x^2+y^2+z^2=4,z\ge0\)，是上半球面。
\emph{条件与易错：}等号处函数值为 \(0\)，不能删边界；定义域是圆盘而非球。
\emph{核验：}\((2,0)\) 合法而 \((2,1)\) 不合法。

\textbf{例2（开边界）。}
\emph{信号：}\(f(x,y)=\ln(y-x^2)\)。
\emph{分析与解答：}对数要求 \(y-x^2>0\)，故定义域是抛物线
\(y=x^2\) 上方且不含抛物线的区域。
\emph{条件与易错：}严格大于号决定边界不属于定义域；不能只写 \(y>0\)。
\emph{核验：}\((0,1)\) 合法，而 \((1,1)\) 使真数为 \(0\)。

\textbf{例3（约束取交）。}
\emph{信号：}
\[
f(x,y)=\frac{\arcsin(x+y)}{\sqrt{\,4-x^2-y^2\,}}.
\]
\emph{分析与解答：}\(\arcsin\) 要求 \(-1\le x+y\le1\)，分母根式要求
\(x^2+y^2<4\)。定义域是开圆盘与闭条带的交集。
\emph{条件与易错：}分母使圆周必须删除，而条带两条边界可以保留。
\emph{核验：}\((1,0)\) 合法，\((2,0)\) 因分母为零不合法。

\textbf{例4（由等值线理解图形）。}
\emph{信号：}\(z=x^2+y^2\)。
\emph{分析与解答：}定义域为 \(\mathbb R^2\)，固定 \(z=c\ge0\) 得
\(x^2+y^2=c\)，水平截线为圆，因此图形是开口向上的旋转抛物面。
\emph{条件与易错：}\(c<0\) 时无截线；不能把等值线圆误认成整个图形。
\emph{核验：}沿任意过原点直线 \(y=kx\)，截得 \(z=(1+k^2)x^2\)，均为抛物线。

\textbf{失分防线、考场表达与掌握标准。}
错误写法“\(x,y\in[-2,2]\)”把圆盘扩大为正方形；快速检查应代入一个角点。
考场写法固定为“由根号、对数、分母分别得……，取交得
\(D=\{\cdots\}\)，边界……”。基础过关能列全约束；130 分以上能准确画区域并
标边界；满分能力能由截痕或等值线解释三维图形并反查定义域。
\textbf{连接：}本节的距离与区域语言直接用于
\knowledgeRef[MATH1-KN-CALC-0180]{二重极限定义}。
\end{knowledgeBox}

\studySubsection{calc-13-k133}{K133 二重极限定义}
\knowledgeAnchor[MATH1-KN-CALC-0180]{二重极限定义}

\begin{knowledgeBox}
\textbf{定位与前置桥。}
一元极限控制的是 \(|x-x_0|\)，二重极限控制点到点的距离
\[
\rho=\sqrt{(x-x_0)^2+(y-y_0)^2}.
\]
距离是什么：它衡量同时改变两个变量后的总偏移；为何需要：趋近可沿无穷多条
路径发生；怎样使用：把误差压成只依赖 \(\rho\) 的量；最小模型：
\(|x-x_0|\le\rho,|y-y_0|\le\rho\)。

\textbf{严格定义。}
\(\lim_{(x,y)\to(x_0,y_0)}f(x,y)=A\) 是指对任意
\(\varepsilon>0\)，存在 \(\delta>0\)，当定义域内
\(0<\rho<\delta\) 时都有 \(|f(x,y)-A|<\varepsilon\)。
“任意路径”已经包含在这个统一的 \(\delta\) 中；有限条路径只能提供必要检查。

\textbf{例1（直接按距离）。}
\emph{信号：}\(f=x^2+y^2\)。
\emph{分析与解答：}在原点有 \(|f|=\rho^2\)。给定 \(\varepsilon>0\)，取
\(\delta=\sqrt{\varepsilon}\)，则 \(0<\rho<\delta\Rightarrow|f|<\varepsilon\)。
\emph{条件与易错：}\(\delta\) 只依赖 \(\varepsilon\)，不能依赖路径。
\emph{核验：}极坐标下 \(f=r^2\to0\)。

\textbf{例2（夹逼估计）。}
\emph{信号：}
\(\displaystyle f(x,y)=\frac{xy}{\sqrt{x^2+y^2}}\)。
\emph{分析与解答：}由 \(2|xy|\le x^2+y^2=\rho^2\)，
\[
|f|\le\frac{\rho^2}{2\rho}=\frac{\rho}{2}\to0.
\]
\emph{条件与易错：}只讨论 \((x,y)\ne(0,0)\) 的去心邻域；分母不能先约掉。
\emph{核验：}沿 \(y=x\) 得 \(|x|/\sqrt2\to0\)。

\textbf{例3（迭代极限不代替二重极限）。}
\emph{信号：}\(\displaystyle f=\frac{xy}{x^2+y^2}\)。
\emph{分析与解答：}先令 \(x\to0\) 再令 \(y\to0\) 得 \(0\)，反序也得 \(0\)；
但沿 \(y=x\) 恒为 \(1/2\)，故二重极限不存在。
\emph{条件与易错：}两个迭代极限相等仍不足以证明二重极限存在。
\emph{核验：}沿 \(y=0\) 值为 \(0\)，已经与 \(y=x\) 冲突。

\textbf{例4（非原点平移）。}
\emph{信号：}\(\lim_{(x,y)\to(1,-1)}(x^2+xy)\)。
\emph{分析与解答：}多项式在该点连续，可代入得 \(1-1=0\)；若按定义，
令 \(u=x-1,v=y+1\)，误差可估为
\(|u-v+u^2+uv|\le2\rho+2\rho^2\to0\)。
\emph{条件与易错：}平移后必须把 \(x=1+u,y=-1+v\) 全部代入。
\emph{核验：}沿任意直线 \(v=ku\) 也趋于 \(0\)。

\textbf{失分防线与考场标准。}
错误写法“沿所有直线都相同，所以极限存在”遗漏曲线路径；正确判据是统一估计。
考场证明写“令 \(\rho=\cdots\)，则 \(|f-A|\le\phi(\rho)\to0\)”。
基础过关能区分二重与迭代极限；130 分以上能造统一界；满分能力能在
\(\varepsilon\)-\(\delta\)、夹逼和反例路径之间快速选法。
\textbf{连接：}否定极限由
\knowledgeRef[MATH1-KN-CALC-0181]{路径法}完成，证明存在常用
\knowledgeRef[MATH1-KN-CALC-0182]{极坐标估计}。
\end{knowledgeBox}

\studySubsection{calc-13-k134}{K134 路径法否定二重极限}
\knowledgeAnchor[MATH1-KN-CALC-0181]{路径法否定二重极限}

\begin{knowledgeBox}
\textbf{前置短桥与核心。}
\knowledgeRef[MATH1-KN-CALC-0180]{二重极限定义}要求所有合法趋近方式
得到同一个值；它的作用是给路径法提供反证依据。
识别方式是看分子分母的齐次次数或某个组合 \(y/x^p\)；最小反例
\(\frac{xy}{x^2+y^2}\) 沿 \(y=0\) 与 \(y=x\) 取不同极限。
因此，只要找到两条合法路径极限不同，或一条路径上的极限不存在，就可否定；
有限条路径相同永远不能正面证明存在。

\textbf{方法选择：}只需否定时先试两条直线；若齐次结构暗示
\(y/x^p\)，再构造 \(y=kx^p\) 曲线族。若目标是证明极限存在，
就停止枚举路径，转用 K135 的统一估计。

\textbf{例1（直线族）。}
\emph{信号：}\(\displaystyle f=\frac{xy}{x^2+y^2}\)。
\emph{分析与解答：}沿 \(y=kx\) 有
\[
f(x,kx)=\frac{k}{1+k^2},
\]
结果依赖 \(k\)，故原极限不存在。
\emph{条件与易错：}\(x\ne0\) 时路径合法；不能把某个固定 \(k\) 的结果当总极限。
\emph{核验：}\(k=0\) 得 \(0\)，\(k=1\) 得 \(1/2\)。

\textbf{例2（抛物线路径）。}
\emph{信号：}\(\displaystyle f=\frac{x^2y}{x^4+y^2}\)。
\emph{分析与解答：}沿所有直线 \(y=kx\) 都趋于 \(0\)，但沿
\(y=kx^2\) 得 \(k/(1+k^2)\)，依赖 \(k\)，故极限不存在。
\emph{条件与易错：}“直线全过关”不是证明；次数提示应令 \(y\) 与 \(x^2\) 同阶。
\emph{核验：}\(k=0,1\) 分别给 \(0,1/2\)。

\textbf{例3（比例路径）。}
\emph{信号：}\(\displaystyle f=\frac{x-y}{|x|+|y|}\)。
\emph{分析与解答：}沿 \(y=0,x\to0^+\) 得 \(1\)，沿
\(x=0,y\to0^+\) 得 \(-1\)，故极限不存在。
\emph{条件与易错：}绝对值使左右方向也可能不同，要写清单侧参数。
\emph{核验：}函数绝对值不超过 \(1\)，但有界不代表有极限。

\textbf{例4（定义域约束下选路）。}
\emph{信号：}在 \(y\ge0\) 内考察
\(\displaystyle f=\frac{y}{\sqrt{x^2+y^2}}\) 于原点的极限。
\emph{分析与解答：}沿合法路径 \(y=0\) 得 \(0\)，沿 \(x=0,y\to0^+\) 得 \(1\)，
故相对该定义域的极限仍不存在。
\emph{条件与易错：}不能选定义域外的 \(y<0\)；边界路径同样属于允许路径。
\emph{核验：}极坐标中 \(f=\sin\theta\)，角度依赖直接暴露冲突。

\textbf{失分防线、考场表达与标准。}
错误是只写“取不同路径”而不给路径上极限；考场必须列两条路径、代入并比较。
基础过关会用坐标轴和 \(y=kx\)；130 分以上能按次数选 \(y=kx^p\)；
满分能力能先检查定义域并判断路径法只承担“否定”而非“证明存在”。
\textbf{连接：}正面证明转入
\knowledgeRef[MATH1-KN-CALC-0182]{极坐标统一估计}。
\end{knowledgeBox}

\studySubsection{calc-13-k135}{K135 极坐标估计二重极限}
\knowledgeAnchor[MATH1-KN-CALC-0182]{极坐标估计二重极限}

\begin{knowledgeBox}
\textbf{前置短桥与严格核心。}
\knowledgeRef[MATH1-KN-CALC-0180]{二重极限定义}要求对所有路径统一控制；
\knowledgeRef[MATH1-KN-CALC-0181]{路径法}只能否定而不能证明存在。
为建立统一估计，写极坐标 \(x=r\cos\theta,y=r\sin\theta\)，其中
\(r=\sqrt{x^2+y^2}\) 是到原点的距离。为何需要：它把所有路径统一为
\(r\to0\)，角度 \(\theta\) 收纳路径差异；怎样使用：寻找与
\(\theta\) 无关的界 \(Cr^\alpha\)；最小例子是
\(|x|,|y|\le r\)。只有当常数 \(C\) 对所有角度统一且 \(\alpha>0\) 时，
才由夹逼得到极限。

\textbf{方法选择：}先比较分子、分母的径向次数；角因子统一有界时
直接压成 \(Cr^\alpha\)。若分母可能随角度为零，先检查定义域，
不能强行使用极坐标夹逼。

\textbf{例1（次数估计）。}
\emph{信号：}\(\displaystyle f=\frac{x^2y}{x^2+y^2}\)。
\emph{分析与解答：}\(|x|,|y|\le r\)，故
\(|f|\le r^3/r^2=r\to0\)。
\emph{条件与易错：}分母是 \(r^2\)，估计时不能把它再粗略成零。
\emph{核验：}极坐标表达为 \(r\cos^2\theta\sin\theta\)。

\textbf{例2（利用基本不等式）。}
\emph{信号：}\(\displaystyle f=\frac{x^2y^2}{x^2+y^2}\)。
\emph{分析与解答：}由 \(x^2y^2\le(x^2+y^2)^2/4=r^4/4\)，
\(|f|\le r^2/4\to0\)。
\emph{条件与易错：}不要分别约掉 \(x\) 或 \(y\)，那会漏坐标轴。
\emph{核验：}沿 \(y=x\) 得 \(x^2/2\to0\)。

\textbf{例3（复合小量）。}
\emph{信号：}
\(\displaystyle f=\frac{1-\cos(x^2+y^2)}{x^2+y^2}\)。
\emph{分析与解答：}令 \(u=r^2\)，则
\((1-\cos u)/u\sim u/2\to0\)。
\emph{条件与易错：}这里调用的一元等价无穷小是在 \(u\to0\) 下，且分母非零。
\emph{核验：}用 \(0\le1-\cos u\le u^2/2\) 可得严格夹逼。

\textbf{例4（角度依赖是报警器）。}
\emph{信号：}\(\displaystyle f=\frac{x^2-y^2}{x^2+y^2}\)。
\emph{分析与解答：}极坐标下 \(f=\cos2\theta\)，没有随 \(r\) 衰减，
沿 \(\theta=0,\pi/2\) 分别为 \(1,-1\)，故极限不存在。
\emph{条件与易错：}换成极坐标不自动证明存在；必须消除角度依赖或给统一界。
\emph{核验：}两条坐标轴路径立即复核。

\textbf{失分防线、考场表达与标准。}
错误写法“令 \(r\to0\)，所以式子趋零”没有给出角度一致性。
考场写成“令 \(r=\sqrt{x^2+y^2}\)，有
\(|f-A|\le Cr^\alpha\to0\)，其中 \(C\) 与 \(\theta\) 无关”。
基础过关会代极坐标；130 分以上会用齐次阶数和基本不等式；
满分能力能识别何时角度项揭示不存在并切回路径法。
\textbf{连接：}统一极限是连续与可微余项检验的共同工具。
\end{knowledgeBox}

\studySubsection{calc-13-k136}{K136 多元连续性与有界闭区域性质}
\knowledgeAnchor[MATH1-KN-CALC-0183]{多元连续性与有界闭区域性质}

\begin{knowledgeBox}
\textbf{前置短桥。}
\knowledgeRef[MATH1-KN-CALC-0180]{二重极限定义}给出“任意路径统一趋近”的含义；
连续要求这个极限恰等于函数在该点的值。
它是什么：局部输入小变化引起输出小变化；为何需要：合法代入、复合运算和最值存在；
怎样识别：初等连续函数在其定义域内按四则与复合保持连续；最小反例是把原点函数值
改成与极限不同的常数。

\textbf{严格核心。}
\(f\) 在 \(P\) 连续指 \(\lim_{Q\to P,Q\in D}f(Q)=f(P)\)。
若 \(D\) 是有界闭区域且 \(f\) 在 \(D\) 上连续，则 \(f\) 在 \(D\) 上有界并取得
最大值、最小值。该定理只保证存在，不自动给出数值；缺“闭”“有界”或“连续”
任一条件都可能失败。

\textbf{方法选择：}初等连续函数的四则与复合优先直接代入；分段点
必须核验“极限等于函数值”；闭区域最值题先逐项确认闭、有界、连续，
再进入候选点比较。

\textbf{例1（分段点连续）。}
\emph{信号：}
\[
f(x,y)=
\begin{cases}
\dfrac{x^2y}{x^2+y^2},&(x,y)\ne(0,0),\\
0,&(x,y)=(0,0).
\end{cases}
\]
\emph{分析与解答：}\(|f|\le |y|\le r\to0=f(0,0)\)，故在原点连续。
\emph{条件与易错：}分母只在去心邻域使用；沿几条路径为零不足以完成证明。
\emph{核验：}统一界不依赖路径。

\textbf{例2（复合连续）。}
\emph{信号：}\(f(x,y)=\ln(1+x^2+y^2)\)。
\emph{分析与解答：}\(1+x^2+y^2>0\) 对全平面成立，多项式与对数复合连续，
故 \(f\) 在 \(\mathbb R^2\) 连续。
\emph{条件与易错：}必须先检查对数定义域，再使用复合连续性。
\emph{核验：}任一点附近真数均至少为 \(1\)。

\textbf{例3（只证最值存在）。}
\emph{信号：} \(f=x^2-xy+y^2\) 在
\(D=\{x^2+y^2\le1\}\) 上。
\emph{分析与解答：}多项式连续，\(D\) 闭且有界，故 \(f\) 在 \(D\) 上必取得
最大、最小值。
\emph{条件与易错：}本问若只要求存在性，无须提前求驻点；“有界”不等于“取到”。
\emph{核验：}原点属于 \(D\) 且 \(f(0,0)=0\)，至少给出一个候选值。

\textbf{例4（条件缺失的反例）。}
\emph{信号：} \(f(x,y)=x\) 在开圆盘 \(x^2+y^2<1\) 上。
\emph{分析与解答：}函数连续且有界，但上确界为 \(1\)，圆盘内没有 \(x=1\) 的点，
故最大值不存在。
\emph{条件与易错：}区域不闭导致边界候选丢失；不能把上确界写成最大值。
\emph{核验：}点 \((1-1/n,0)\) 的函数值趋于 \(1\)。

\textbf{失分防线、考场表达与标准。}
错误是只写“连续所以有最值”，漏掉区域闭且有界。
考场模板：“\(f\) 在有界闭区域 \(D\) 上连续，由最值定理，最大值、最小值均存在。”
基础过关会判点连续；130 分以上会给统一估计并核查区域；
满分能力能区分最大值与上确界、存在性与求值任务。
\textbf{连接：}连续是可微的必要后果，也是偏导连续推出可微时的关键条件。
\end{knowledgeBox}

\studySubsection{calc-13-k137}{K137 一阶偏导数定义与计算}
\knowledgeAnchor[MATH1-KN-CALC-0184]{一阶偏导数定义与计算}

\begin{knowledgeBox}
\textbf{定位与前置桥。}
一元导数测沿数轴的瞬时变化；偏导数是在多元函数中固定其余变量，只沿一个
坐标方向求一元导数。它是什么：\(f_x\) 是 \(y\) 固定后的 \(x\) 方向变化率；
为何需要：它组成全微分、梯度和切平面；怎样使用：普通点用求导公式，分段点回定义；
最小例子 \(f=x^2+y\) 有 \(f_x=2x,f_y=1\)。

\textbf{严格定义与条件。}
\[
f_x(x_0,y_0)=\lim_{h\to0}
\frac{f(x_0+h,y_0)-f(x_0,y_0)}h,
\quad
f_y(x_0,y_0)=\lim_{h\to0}
\frac{f(x_0,y_0+h)-f(x_0,y_0)}h.
\]
偏导只考察两条坐标方向，偏导存在不保证连续或可微。

\textbf{方法选择：}普通点直接把其余变量视为常数求导；分段点和特殊点
回到差商定义；抽象复合结构先写清哪些变量在本次偏导中固定。

\textbf{例1（直接计算）。}
\emph{信号：}\(f=x^2y+e^{xy}\)。
\emph{分析与解答：}固定 \(y\) 得
\(f_x=2xy+ye^{xy}\)；固定 \(x\) 得
\(f_y=x^2+xe^{xy}\)。
\emph{条件与易错：}求 \(f_x\) 时 \(y\) 是常数，但 \(e^{xy}\) 仍需链式求导。
\emph{核验：}在 \((0,0)\) 两式均为 \(0\)。

\textbf{例2（特殊点回定义）。}
\emph{信号：}\(f(x,y)=|x|+y^2\)，求原点偏导。
\emph{分析与解答：}
\[
\frac{f(h,0)-f(0,0)}h=\frac{|h|}{h}
\]
左右极限不同，故 \(f_x(0,0)\) 不存在；而
\([f(0,h)-f(0,0)]/h=h\to0\)，故 \(f_y(0,0)=0\)。
\emph{条件与易错：}不可在 \(|x|\) 的尖点直接套符号函数而忽略双侧极限。
\emph{核验：}图形沿 \(x\) 轴有尖点，沿 \(y\) 轴为光滑抛物线。

\textbf{例3（偏导存在但函数不连续）。}
\emph{信号：}
\[
f(x,y)=
\begin{cases}\dfrac{xy}{x^2+y^2},&(x,y)\ne(0,0),\\0,&(0,0).\end{cases}
\]
\emph{分析与解答：}沿两坐标轴函数恒为 \(0\)，故
\(f_x(0,0)=f_y(0,0)=0\)；但沿 \(y=x\) 函数为 \(1/2\)，在原点不连续。
\emph{条件与易错：}两个偏导只看两条路径，不能推出二重极限。
\emph{核验：}坐标轴差商均确为零。

\textbf{例4（先改写再求偏导）。}
\emph{信号：}\(f(x,y)=x^y\)，其中 \(x>0\)。
\emph{分析与解答：}写成 \(e^{y\ln x}\)，故
\[
f_x=yx^{y-1},\qquad f_y=x^y\ln x.
\]
\emph{条件与易错：}实变量下为使用 \(\ln x\) 必须写 \(x>0\)；不能把 \(y\)
当常数后误认为 \(f_y=0\)。
\emph{核验：}取 \(y=2\) 时 \(f_x=2x\) 与一元结果一致。

\textbf{失分防线、考场表达与标准。}
错误是求 \(f_x\) 时让 \(y\) 随 \(x\) 一起变化；考场可写“固定 \(y\)，由定义/公式”。
基础过关能准确固定变量；130 分以上会在分段点回定义；
满分能力能主动说明偏导存在与连续、可微没有充分关系。
\textbf{连接：}偏导是
\knowledgeRef[MATH1-KN-CALC-0186]{全微分}和
\knowledgeRef[MATH1-KN-CALC-0196]{梯度}的分量。
\end{knowledgeBox}

\studySubsection{calc-13-k138}{K138 高阶偏导与混合偏导}
\knowledgeAnchor[MATH1-KN-CALC-0185]{高阶偏导与混合偏导}

\begin{knowledgeBox}
\textbf{前置短桥与核心。}
\knowledgeRef[MATH1-KN-CALC-0184]{一阶偏导}先固定一个变量求导；
对所得函数再次偏导就得到
\(f_{xx},f_{xy},f_{yx},f_{yy}\)。下标按求导先后从左向右读：
\(f_{xy}=(f_x)_y\)。混合偏导在点的某邻域内连续是
\(f_{xy}=f_{yx}\) 的常用充分条件，不能无条件交换。

\textbf{方法选择：}计算时按下标顺序逐次求导；要交换混合偏导顺序，
必须先给出连续性等充分条件；命题判断若条件不足，优先寻找不相等反例。

\textbf{例1（四个二阶偏导）。}
\emph{信号：}\(f=e^{xy}\)。
\emph{分析与解答：}
\[
f_{xx}=y^2e^{xy},\quad
f_{xy}=(1+xy)e^{xy},\quad
f_{yx}=(1+xy)e^{xy},\quad
f_{yy}=x^2e^{xy}.
\]
\emph{条件与易错：}计算 \(f_{xy}\) 时要同时求导系数 \(y\) 与指数。
\emph{核验：}函数光滑，两个混合偏导应相等。

\textbf{例2（读清下标顺序）。}
\emph{信号：}\(f=x^3y^2+xy\)。
\emph{分析与解答：}
\(f_x=3x^2y^2+y\)，再对 \(y\) 求导得
\(f_{xy}=6x^2y+1\)；反序同样得到 \(6x^2y+1\)。
\emph{条件与易错：}不能把 \(f_{xy}\) 读成先 \(y\) 后 \(x\) 后仍沿用旧中间式。
\emph{核验：}多项式的各阶偏导连续。

\textbf{例3（混合偏导可以不相等）。}
\emph{信号：}
\[
f(x,y)=
\begin{cases}
\dfrac{xy(x^2-y^2)}{x^2+y^2},&(x,y)\ne(0,0),\\0,&(0,0).
\end{cases}
\]
\emph{分析与解答：}先算
\(f_x(0,y)=-y\)，故 \(f_{xy}(0,0)=-1\)；
又 \(f_y(x,0)=x\)，故 \(f_{yx}(0,0)=1\)。
\emph{条件与易错：}原点附近缺少保证混合偏导连续的条件，不能先交换再算。
\emph{核验：}两个定义差商给出相反符号，直接证明不相等。

\textbf{例4（由混合偏导相等求参数）。}
\emph{信号：}已知
\(f_x=ay+2xy,\ f_y=3x+x^2\)，问何时可能来自同一 \(C^2\) 函数。
\emph{分析与解答：}
\((f_x)_y=a+2x\)，\((f_y)_x=3+2x\)，故必须且在此题中可取
\(a=3\)。
\emph{条件与易错：}相等是局部相容条件；题设 \(C^2\) 才允许直接使用。
\emph{核验：}当 \(a=3\) 时可积分得 \(f=3xy+x^2y+C\)。

\textbf{失分防线、考场表达与标准。}
错误写法是“混合偏导必相等”；正确写法要附“在邻域连续”等充分条件。
基础过关会算四个二阶偏导；130 分以上能用相容性求参数；
满分能力能在特殊点按定义构造不相等反例。
\textbf{连接：}高阶偏导进入二阶复合求导与多元 Taylor 工具。
\end{knowledgeBox}

\studySubsection{calc-13-k139}{K139 全微分定义}
\knowledgeAnchor[MATH1-KN-CALC-0186]{全微分定义}

\begin{knowledgeBox}
\textbf{定位与前置桥。}
\knowledgeRef[MATH1-KN-CALC-0184]{偏导}只描述坐标方向；
全微分要求用同一个线性式同时逼近所有方向的函数增量。
令 \(\Delta x,\Delta y\) 为输入增量，
\(\rho=\sqrt{\Delta x^2+\Delta y^2}\) 为总距离。\(o(\rho)\) 的意思是余项除以
\(\rho\) 后趋于零；最小例子 \(f=x^2+y^2\) 在原点的增量为
\(\rho^2=o(\rho)\)。

\textbf{严格定义。}
若
\[
\Delta z=A\Delta x+B\Delta y+o(\rho),
\]
则 \(f\) 在该点可微，线性主部
\(dz=A\,dx+B\,dy\) 称全微分；此时必有
\(A=f_x,B=f_y\)。余项必须对所有方向统一为高阶，不能只沿坐标轴检验。

\textbf{方法选择：}光滑显式函数直接用偏导写线性主部；特殊点则把
\(\Delta z-f_x\Delta x-f_y\Delta y\) 除以 \(\rho\) 作统一估计。
线性近似题最后还要说明余项是 \(o(\rho)\)。

\textbf{例1（多项式的全微分）。}
\emph{信号：}\(f=x^2+xy\)，点 \((1,2)\)。
\emph{分析与解答：}\(f_x=2x+y,f_y=x\)，故
\[
dz=4\,dx+1\,dy.
\]
\emph{条件与易错：}系数必须在指定点求值。
\emph{核验：}直接展开增量得
\(\Delta z=4\Delta x+\Delta y+\Delta x^2+\Delta x\Delta y\)，余项为 \(O(\rho^2)\)。

\textbf{例2（线性近似）。}
\emph{信号：}\(f=\sqrt{4-x^2-y^2}\)，在 \((1,1)\) 附近估值。
\emph{分析与解答：}该点 \(f=\sqrt2\)，
\(f_x=f_y=-1/\sqrt2\)，故
\[
f(1+\Delta x,1+\Delta y)
\approx\sqrt2-\frac{\Delta x+\Delta y}{\sqrt2}.
\]
\emph{条件与易错：}邻域需保持根号内为正；近似号不是恒等号。
\emph{核验：}取小增量代回原式可检查误差为二阶。

\textbf{例3（原点的零线性主部）。}
\emph{信号：}\(f=x^2+y^2\) 于原点。
\emph{分析与解答：}\(f_x=f_y=0\)，且
\(\Delta z=\rho^2=o(\rho)\)，故可微且 \(dz=0\)。
\emph{条件与易错：}全微分为零不表示函数在邻域恒定。
\emph{核验：}\(\Delta z/\rho=\rho\to0\)。

\textbf{例4（误差传播）。}
\emph{信号：}矩形面积 \(S=xy\)，测量误差为 \(dx,dy\)。
\emph{分析与解答：}
\[
dS=y\,dx+x\,dy,
\]
它给出一阶绝对误差；精确增量还多 \(dx\,dy\)，该项相对
\(\rho\) 为高阶。
\emph{条件与易错：}若要求严格误差界，不能把 \(dS\) 当精确增量。
\emph{核验：}\((x+dx)(y+dy)-xy\) 直接展开。

\textbf{失分防线、考场表达与标准。}
错误是写成“分别对 \(dx,dy\) 高阶”，却没有统一距离。
考场定义必须出现线性主部和 \(o(\rho)\)。
基础过关会求 \(dz\)；130 分以上能由增量验证；
满分能力能解释可微为何是全方向统一线性近似。
\textbf{连接：}逻辑关系由
\knowledgeRef[MATH1-KN-CALC-0187]{可微、连续与偏导关系}统一。
\end{knowledgeBox}

\studySubsection{calc-13-k140}{K140 可微、连续、偏导存在的逻辑关系}
\knowledgeAnchor[MATH1-KN-CALC-0187]{可微、连续、偏导存在的逻辑关系}

\begin{knowledgeBox}
\textbf{前置短桥与逻辑图。}
\knowledgeRef[MATH1-KN-CALC-0184]{偏导}是坐标方向差商，
\knowledgeRef[MATH1-KN-CALC-0183]{连续}是函数值的统一极限，
\knowledgeRef[MATH1-KN-CALC-0186]{全微分定义}把可微表述为全方向统一线性近似。
三者作用不同。必须掌握
\[
\text{{可微}}\Rightarrow\text{{连续且偏导存在}},
\qquad
\text{{邻域内偏导存在并在点连续}}\Rightarrow\text{{可微}}.
\]
第二条是充分条件而非必要条件；“偏导存在”“连续”单独都不能推出可微。

\textbf{方法选择：}逻辑判断只沿已证明的箭头推；要否定逆命题就选
连续性、偏导或可微性相互分离的反例；要正面证明可微，可用偏导连续
充分条件，条件不足时转 K141 的定义检验。

\textbf{例1（充分条件的标准使用）。}
\emph{信号：}\(f=x^2y+\sin(xy)\)。
\emph{分析与解答：}两个一阶偏导由初等连续函数组成，在全平面连续，故 \(f\)
处处可微，进而处处连续。
\emph{条件与易错：}要说偏导在邻域存在并在点连续，而非只在该点有数值。
\emph{核验：}函数本身也是初等连续函数。

\textbf{例2（偏导存在但不连续）。}
\emph{信号：}\(f=xy/(x^2+y^2)\) 在原点补零。
\emph{分析与解答：}两偏导在原点均为 \(0\)，但沿 \(y=x\) 函数为 \(1/2\)，
不连续，因而必不可微。
\emph{条件与易错：}不可由两条坐标轴的信息替代全方向。
\emph{核验：}利用“可微必连续”可立即否定。

\textbf{例3（连续且偏导存在仍可不可微）。}
\emph{信号：}
\[
f(x,y)=
\begin{cases}\dfrac{x^3}{x^2+y^2},&(x,y)\ne(0,0),\\0,&(0,0).\end{cases}
\]
\emph{分析与解答：}\(|f|\le|x|\le r\)，故连续；
\(f_x(0,0)=1,f_y(0,0)=0\)。候选线性主部为 \(x\)，沿 \(y=x\) 有
\[
\frac{f-x}{\sqrt{x^2+y^2}}
=\frac{-x/2}{\sqrt2|x|}
\]
不趋零，故不可微。
\emph{条件与易错：}连续加偏导存在仍缺统一线性余项。
\emph{核验：}定义检验给出非零常量。

\textbf{例4（可微但偏导不连续）。}
\emph{信号：}令 \(r=\sqrt{x^2+y^2}\)，
\[
f(x,y)=
\begin{cases}r^2\sin(1/r),&r>0,\\0,&r=0.\end{cases}
\]
\emph{分析与解答：}\(|f|/r\le r\to0\)，故原点可微且全微分为零；
但非原点偏导含 \(\cos(1/r)\) 型振荡项，不能在原点连续。
\emph{条件与易错：}偏导连续只是方便的充分条件，不是可微的必要条件。
\emph{核验：}可微性直接由余项定义完成，不依赖偏导连续。

\textbf{失分防线、考场表达与标准。}
禁止写双向等价箭头。考场逻辑题先画单向图，再用标准反例拆逆命题。
基础过关记对蕴含；130 分以上能选反例；
满分能力能分别用连续性、偏导定义和余项定义给出证据。
\textbf{连接：}遇到特殊分段函数时进入
\knowledgeRef[MATH1-KN-CALC-0188]{可微性的定义检验}。
\end{knowledgeBox}

\studySubsection{calc-13-k141}{K141 可微性的定义检验}
\knowledgeAnchor[MATH1-KN-CALC-0188]{可微性的定义检验}

\begin{knowledgeBox}
\textbf{前置短桥与算法。}
\knowledgeRef[MATH1-KN-CALC-0186]{全微分定义}要求函数增量减去候选线性主部后
为 \(o(\rho)\)；\knowledgeRef[MATH1-KN-CALC-0187]{三者逻辑关系}
只给必要条件，不能代替下面的余项检验。
候选系数若存在只能是偏导，因此标准流程是：
先查连续性，按定义求偏导，再检验
\[
\frac{\Delta f-f_x(P)\Delta x-f_y(P)\Delta y}{\rho}\to0.
\]
一条路径给出非零极限可否定；正面证明必须统一估计。

\textbf{例1（连续但不可微）。}
\emph{信号：}
\[
f(x,y)=
\begin{cases}\dfrac{x^2y}{x^2+y^2},&(x,y)\ne(0,0),\\0,&(0,0).\end{cases}
\]
\emph{分析与解答：}\(|f|\le|y|\le r\)，函数连续，且两偏导均为 \(0\)。
沿 \(y=x\)，\(f=x/2\)，于是 \(f/r\) 的绝对值趋于 \(1/(2\sqrt2)\)，
故不可微。
\emph{条件与易错：}连续和偏导为零只完成必要条件检查。
\emph{核验：}路径给出余项检验失败的确定常数。

\textbf{例2（零线性主部下可微）。}
\emph{信号：}\(f=(x^2+y^2)^{3/4}=r^{3/2}\)。
\emph{分析与解答：}原点两偏导为 \(0\)，且
\[
\frac{|f|}{r}=r^{1/2}\to0,
\]
故原点可微，全微分为零。
\emph{条件与易错：}函数虽含分数次幂，但 \(r\ge0\) 且定义良好。
\emph{核验：}统一表达只依赖 \(r\)。

\textbf{例3（非零线性主部）。}
\emph{信号：}
\[
f(x,y)=2x-y+
\begin{cases}(x^2+y^2)\sin\dfrac1{\sqrt{x^2+y^2}},&(x,y)\ne(0,0),\\0,&(0,0).\end{cases}
\]
\emph{分析与解答：}候选线性主部是 \(2x-y\)，余项绝对值不超过 \(r^2\)，
除以 \(r\) 后趋零，故可微且 \(df=2dx-dy\)。
\emph{条件与易错：}必须先减掉线性项，不能直接拿 \(f/r\) 检验。
\emph{核验：}余项统一界 \(r^2\) 与方向无关。

\textbf{例4（参数阈值）。}
\emph{信号：}\(f_\alpha=(x^2+y^2)^\alpha=r^{2\alpha}\)，原点补零。
\emph{分析与解答：}当 \(\alpha>0\) 时两偏导在原点是否为 \(0\) 需进一步看轴上
\(|h|^{2\alpha-1}\)；偏导存在要求 \(\alpha>1/2\)，
而可微检验为 \(r^{2\alpha-1}\to0\)，恰要求 \(\alpha>1/2\)。
\emph{条件与易错：}\(\alpha=1/2\) 时 \(f=r\)，轴上偏导已不存在；
\(\alpha<1/2\) 更不可能可微。
\emph{核验：}径向表达直接给出充分必要阈值。

\textbf{失分防线、考场表达与标准。}
错误是只写“连续且偏导存在，所以可微”。考场模板应完整写候选线性主部和余项商。
基础过关会按三步检验；130 分以上能用路径否定或统一估计证明；
满分能力能处理非零线性主部与参数阈值。
\textbf{连接：}可微建立后才可安全调用多元链式法则、方向导数公式和切平面。
\end{knowledgeBox}

\studySubsection{calc-13-k142}{K142 一阶链式法则}
\knowledgeAnchor[MATH1-KN-CALC-0189]{一阶链式法则}

\begin{knowledgeBox}
\textbf{定位与前置桥。}
一元链式法则把“外层变化率”乘“内层变化率”；多元情形要把从自变量到结果的
每一条依赖路径相加。偏导是什么、为何需要与最小模型已见
\knowledgeRef[MATH1-KN-CALC-0184]{K137}；本节只补作用：偏导是每条路径上的局部倍率。
若 \(z=f(u,v)\)，且 \(u=u(x,y),v=v(x,y)\)，各函数在相应点可微，则
\[
z_x=f_u u_x+f_v v_x,\qquad z_y=f_u u_y+f_v v_y.
\]
所有偏导都必须在对应的 \((u,v)\) 与 \((x,y)\) 点求值。

\textbf{方法选择：}只有一个自变量时写全导数；有多个自变量时画变量树，
把每条依赖路径上的偏导乘积相加；抽象外层函数保留 \(f_u,f_v\)，
不要擅自把它们当常数。

\textbf{例1（抽象函数与变量树）。}
\emph{信号：}\(z=f(u,v),u=x+y,v=xy\)。
\emph{分析与解答：}
\[
z_x=f_u+yf_v,\qquad z_y=f_u+xf_v,
\]
其中 \(f_u,f_v\) 在 \((x+y,xy)\) 取值。
\emph{条件与易错：}从 \(x\) 到 \(z\) 有经 \(u,v\) 两条路径，不能漏 \(yf_v\)。
\emph{核验：}取 \(f=u+v\) 时直接求导与上式一致。

\textbf{例2（一个中间变量）。}
\emph{信号：}\(z=e^u,u=x^2+y^2\)。
\emph{分析与解答：}
\[
z_x=2xe^{x^2+y^2},\qquad z_y=2ye^{x^2+y^2}.
\]
\emph{条件与易错：}外层导数 \(e^u\) 仍要把 \(u\) 回代。
\emph{核验：}直接视为 \(e^{x^2+y^2}\) 求偏导。

\textbf{例3（多个独立变量）。}
\emph{信号：}\(w=f(xy,yz,zx)\)，求 \(w_x\)。
\emph{分析与解答：}令 \(u=xy,v=yz,s=zx\)，则
\[
w_x=f_u y+f_v\cdot0+f_s z=yf_u+zf_s.
\]
\emph{条件与易错：}\(v=yz\) 不依赖 \(x\)，该路径贡献为零；三个外层偏导的求值点
是 \((xy,yz,zx)\)。
\emph{核验：}取 \(f=u+v+s\)，直接导数为 \(y+z\)。

\textbf{例4（由链式法则验证关系）。}
\emph{信号：}\(z=f(x+y,x-y)\)。
\emph{分析与解答：}令 \(u=x+y,v=x-y\)，得
\[
z_x=f_u+f_v,\qquad z_y=f_u-f_v,
\]
故 \(z_x+z_y=2f_u\)，\(z_x-z_y=2f_v\)。
\emph{条件与易错：}两个 \(v\) 路径的符号相反；不能把 \(v_y\) 写成 \(1\)。
\emph{核验：}用 \(f=u^2+v^2\) 可直接复核。

\textbf{失分防线、考场表达与标准。}
错误是凭记忆写公式而漏路径或漏求值点。考场先画
\(x,y\to u,v\to z\) 变量树，再沿路径乘、对路径求和。
基础过关会处理一个中间变量；130 分以上能处理抽象多路径；
满分能力能从链式结果反推偏微分关系并核对求值点。
\textbf{连接：}再次对一阶结果求导进入
\knowledgeRef[MATH1-KN-CALC-0190]{二阶复合偏导}。
\end{knowledgeBox}

\studySubsection{calc-13-k143}{K143 二阶复合偏导}
\knowledgeAnchor[MATH1-KN-CALC-0190]{二阶复合偏导}

\begin{knowledgeBox}
\textbf{前置短桥与算法。}
\knowledgeRef[MATH1-KN-CALC-0189]{一阶链式法则}沿每条依赖路径相乘并求和；
\knowledgeRef[MATH1-KN-CALC-0185]{高阶偏导}要求对所得一阶式再次逐项求导；
二阶偏导不是简单把一阶系数平方，
而是对一阶结果逐项再求导。这样既产生外层二阶偏导与两条一阶路径的乘积，
也产生外层一阶偏导乘内层二阶导。最稳流程是“先写一阶，保留中间变量，再逐项求导”。

\textbf{例1（一个非线性中间变量）。}
\emph{信号：}\(z=f(u),u=x^2+y^2\)。
\emph{分析与解答：}
\[
z_x=2xf'(u),\qquad
z_{xx}=2f'(u)+4x^2f''(u).
\]
\emph{条件与易错：}第一项来自 \(u_{xx}=2\)，只平方 \(u_x\) 会漏掉它。
\emph{核验：}取 \(f(u)=u\)，应有 \(z_{xx}=2\)。

\textbf{例2（混合偏导的交叉项）。}
\emph{信号：}\(z=f(u,v),u=x+y,v=xy\)，求 \(z_{xy}\)。
\emph{分析与解答：}从 \(z_x=f_u+yf_v\) 出发，
\[
\begin{aligned}
z_{xy}
&=f_{uu}u_y+f_{uv}v_y+f_v
 +y(f_{vu}u_y+f_{vv}v_y)\\
&=f_{uu}+x f_{uv}+f_v+y f_{vu}+xy f_{vv}.
\end{aligned}
\]
\emph{条件与易错：}求 \(y f_v\) 必须使用乘积法则，不能漏 \(f_v\)。
\emph{核验：}取 \(f=u+v\)，结果为 \(1\)，与 \(z=x+y+xy\) 相符。

\textbf{例3（线性变换的简化）。}
\emph{信号：}\(z=f(x+y,x-y)\)。
\emph{分析与解答：}因中间变量为线性函数，无内层二阶项，
\[
z_{xx}=f_{uu}+2f_{uv}+f_{vv},\qquad
z_{yy}=f_{uu}-2f_{uv}+f_{vv},
\]
故 \(z_{xx}-z_{yy}=4f_{uv}\)（假定混合偏导可交换）。
\emph{条件与易错：}交叉项系数 \(2\) 来自两条次序。
\emph{核验：}取 \(f(u,v)=uv\)，两边均为 \(4\)。

\textbf{例4（抽象函数的二阶径向导数）。}
\emph{信号：}\(z=f(r),r=\sqrt{x^2+y^2}\)，\(r>0\)。
\emph{分析与解答：}
\[
z_x=f'(r)\frac{x}{r},\qquad
z_{xx}=f''(r)\frac{x^2}{r^2}
+f'(r)\frac{y^2}{r^3}.
\]
\emph{条件与易错：}公式只在 \(r>0\) 使用；原点必须另按定义讨论。
\emph{核验：}取 \(f(r)=r^2\)，化简得 \(z_{xx}=2\)。

\textbf{失分防线、考场表达与标准。}
常见错误是漏内层二阶导、漏乘积法则、把交叉项系数写成 \(1\)。
考场先写一阶表达式，再对目标变量逐项求导并在最后回代。
基础过关会处理单中间变量；130 分以上能处理 \(z_{xy}\)；
满分能力能用变量树或矩阵观点完成自检而不漏项。
\textbf{连接：}一阶微分可由
\knowledgeRef[MATH1-KN-CALC-0191]{全微分形式不变性}压缩，但二阶不能照搬。
\end{knowledgeBox}

\studySubsection{calc-13-k144}{K144 全微分形式不变性}
\knowledgeAnchor[MATH1-KN-CALC-0191]{全微分形式不变性}

\begin{knowledgeBox}
\textbf{前置短桥与核心。}
\knowledgeRef[MATH1-KN-CALC-0186]{全微分}是函数增量的一阶线性主部；
\knowledgeRef[MATH1-KN-CALC-0189]{链式法则}把复合关系的所有一阶路径相加。
若 \(z=f(u,v)\)，而 \(u,v\) 又是其他变量的可微函数，则
\[
dz=f_u\,du+f_v\,dv.
\]
把 \(du,dv\) 展开后恰得到对原变量的全微分，这称一阶全微分形式不变。
它要求相关函数可微；非线性换元下二阶微分还会出现内层二阶项，不能简单照搬。

\textbf{例1（线性换元）。}
\emph{信号：}\(u=x+y,v=x-y,z=f(u,v)\)。
\emph{分析与解答：}
\[
du=dx+dy,\quad dv=dx-dy,
\]
\[
dz=(f_u+f_v)dx+(f_u-f_v)dy.
\]
\emph{条件与易错：} \(f_u,f_v\) 在 \((x+y,x-y)\) 取值。
\emph{核验：}系数与 K142 的 \(z_x,z_y\) 完全一致。

\textbf{例2（直接与间接算法一致）。}
\emph{信号：}\(z=u^2v,u=x+y,v=xy\)。
\emph{分析与解答：}
\[
dz=2uv\,du+u^2\,dv
=2(x+y)xy(dx+dy)+(x+y)^2(y\,dx+x\,dy).
\]
\emph{条件与易错：}得到 \(du,dv\) 后仍要代入原变量，除非题目允许保留。
\emph{核验：}直接展开 \(z=(x+y)^2xy\) 求微分得到同式。

\textbf{例3（极坐标的一阶微分）。}
\emph{信号：}\(z=f(x,y)\)，\(x=r\cos\theta,y=r\sin\theta\)。
\emph{分析与解答：}
\[
dx=\cos\theta\,dr-r\sin\theta\,d\theta,\quad
dy=\sin\theta\,dr+r\cos\theta\,d\theta,
\]
代入 \(dz=f_xdx+f_ydy\) 即得到 \(dr,d\theta\) 表达。
\emph{条件与易错：}在 \(r=0\) 极角不唯一，涉及 \(\theta\) 的结论需谨慎。
\emph{核验：}对 \(f=x^2+y^2=r^2\) 应化为 \(dz=2r\,dr\)。

\textbf{例4（二阶形式不保持的反例）。}
\emph{信号：}\(z=u,u=x^2\)。
\emph{分析与解答：}作为 \(u\) 的函数，\(z_{uu}=0\)；但作为 \(x\) 的函数
\(z=x^2\)，有 \(d^2z=2(dx)^2\)。遗漏项来自 \(z_u\,d^2u\)。
\emph{条件与易错：}一阶形式不变不能无条件推广到二阶非线性换元。
\emph{核验：}直接二阶求导给出非零结果。

\textbf{失分防线、考场表达与标准。}
错误是把 \(du,dv\) 当独立常数或把一阶结论套到二阶。
考场可写“由全微分形式不变性，\(dz=\cdots\)，再代 \(du=\cdots\)”。
基础过关会求复合全微分；130 分以上能选择保留或展开中间微分；
满分能力能说明二阶失效的内层二阶来源。
\textbf{连接：}隐函数微分可把方程整体求全微分后解出目标微分。
\end{knowledgeBox}

\studySubsection{calc-13-k145}{K145 一个方程确定一个隐函数}
\knowledgeAnchor[MATH1-KN-CALC-0192]{一个方程确定一个隐函数}

\begin{knowledgeBox}
\textbf{定位与前置桥。}
显函数直接写 \(z=f(x,y)\)；隐式方程 \(F(x,y,z)=0\) 可能只在某点附近唯一确定
\(z=z(x,y)\)。\knowledgeRef[MATH1-KN-CALC-0186]{全微分}统一方程的一阶变化，
\knowledgeRef[MATH1-KN-CALC-0189]{链式法则}说明对隐式关系求导时各路径如何相加。
若 \(F\) 在点附近连续可微、\(F(P)=0\) 且 \(F_z(P)\ne0\)，则局部可确定
\(z=z(x,y)\)，并有
\[
z_x=-\frac{F_x}{F_z},\qquad z_y=-\frac{F_y}{F_z}.
\]
这些 \(F_x,F_y,F_z\) 是把 \(x,y,z\) 暂视为独立变量后的偏导。

\textbf{方法选择：}只求隐函数偏导时直接用分式公式；题目要求说明
局部存在唯一性时，先核验 \(F(P)=0\)、连续可微与 \(F_z(P)\ne0\)。
若目标偏导为零，公式失效，必须另找可解分支。

\textbf{例1（球面局部分支）。}
\emph{信号：}\(F=x^2+y^2+z^2-3=0\)，点 \((1,1,1)\)。
\emph{分析与解答：}\(F_z=2\ne0\)，故附近可定 \(z(x,y)\)，且
\[
z_x=-\frac{2x}{2z}=-1,\qquad z_y=-1
\]
（在该点取值）。
\emph{条件与易错：}球面全局有上下两支，只能说该点附近的一支。
\emph{核验：}上支 \(z=\sqrt{3-x^2-y^2}\) 直接求导一致。

\textbf{例2（指数隐函数）。}
\emph{信号：}\(e^z+xz+y=0\)，在满足方程且 \(e^z+x\ne0\) 的点。
\emph{分析与解答：}
\[
z_x=-\frac{z}{e^z+x},\qquad
z_y=-\frac1{e^z+x}.
\]
\emph{条件与易错：}求 \(F_x\) 时暂把 \(z\) 独立，故 \(F_x=z\)，不能额外乘 \(z_x\)。
\emph{核验：}沿隐函数全导得到
\((e^z+x)z_x+z=0\)。

\textbf{例3（全微分法）。}
\emph{信号：}\(x^2+xy+z^2=1\)。
\emph{分析与解答：}
\[
(2x+y)\,dx+x\,dy+2z\,dz=0,
\]
若 \(z\ne0\)，
\[
dz=-\frac{2x+y}{2z}dx-\frac{x}{2z}dy.
\]
\emph{条件与易错：}除以 \(2z\) 前必须说明 \(z\ne0\)。
\emph{核验：}系数与 \(-F_x/F_z,-F_y/F_z\) 一致。

\textbf{例4（条件失效）。}
\emph{信号：}\(F=z^2-x^2-y^2=0\) 于原点。
\emph{分析与解答：}\(F_z(0,0,0)=0\)，定理不能保证 \(z=z(x,y)\) 可微；
事实上 \(z=\pm\sqrt{x^2+y^2}\) 有两支，且在原点均不可微。
\emph{条件与易错：} \(F_z=0\) 不自动说明绝对不存在隐函数，只说明该定理和公式失效。
\emph{核验：}同一 \((x,y)\ne(0,0)\) 对应两个 \(z\) 值。

\textbf{失分防线、考场表达与标准。}
错误是分子偏导时把 \(z\) 已当函数，造成重复链式；或漏写 \(F_z\ne0\)。
考场先验 \(F(P)=0,F_z(P)\ne0\)，再写公式并代点。
基础过关会套一阶公式；130 分以上会用全微分交叉核验；
满分能力能区分局部存在、全局多分支和退化点。
\textbf{连接：}多个方程同时确定多个函数时需使用 Jacobian 与线性方程组。
\end{knowledgeBox}

\studySubsection{calc-13-k146}{K146 方程组确定隐函数组}
\knowledgeAnchor[MATH1-KN-CALC-0193]{方程组确定隐函数组}

\begin{knowledgeBox}
\textbf{前置短桥与核心。}
\knowledgeRef[MATH1-KN-CALC-0192]{一个方程确定隐函数}依靠非零的目标偏导
解出一个函数；两个方程要局部解出
\(u(x),v(x)\)，需要关于 \((u,v)\) 的 Jacobian
\[
J=\frac{\partial(F,G)}{\partial(u,v)}
=\begin{vmatrix}F_u&F_v\\G_u&G_v\end{vmatrix}\ne0.
\]
对自变量求导后得到关于 \(u',v'\) 的二元线性方程组；解这个系统比强求显式函数更稳。

\textbf{方法选择：}先计算关于待解变量的 Jacobian；非零时再对原方程组
整体求导，解关于各导数的线性方程组。除非题目明确要求，不必先求
\(u(x),v(x)\) 的显式表达式。

\textbf{例1（线性系统核验）。}
\emph{信号：}\(u+v=x,\ u-v=x^2\)。
\emph{分析与解答：}求导得
\[
u'+v'=1,\qquad u'-v'=2x,
\]
故 \(u'=(1+2x)/2,v'=(1-2x)/2\)。
\emph{条件与易错：}系数矩阵行列式为 \(-2\ne0\)。
\emph{核验：}显式解 \(u=(x+x^2)/2,v=(x-x^2)/2\) 一致。

\textbf{例2（非线性方程组）。}
\emph{信号：}
\[
F=u^2+v-x=0,\qquad G=u+v^2-x^2=0.
\]
\emph{分析与解答：}对 \(x\) 求导：
\[
2u\,u'+v'=1,\qquad u'+2v\,v'=2x.
\]
当 \(4uv-1\ne0\) 时线性系统唯一确定 \(u',v'\)；用 Cramer 法可得
\[
u'=\frac{2v-2x}{4uv-1},\qquad
v'=\frac{4ux-1}{4uv-1}.
\]
\emph{条件与易错：}分母就是对应 Jacobian，必须在指定解点求值。
\emph{核验：}代回两条导数方程。

\textbf{例3（两个自变量中的一个偏导）。}
\emph{信号：}\(F(x,y,u,v)=0,G(x,y,u,v)=0\)，求 \(u_x,v_x\)。
\emph{分析与解答：}
\[
\begin{pmatrix}F_u&F_v\\G_u&G_v\end{pmatrix}
\binom{u_x}{v_x}
=-\binom{F_x}{G_x}.
\]
若 \(J\ne0\)，该系统唯一可解；求 \(y\) 偏导时只替换右端为
\(-\binom{F_y}{G_y}\)。
\emph{条件与易错：}系数矩阵不随所求自变量改变，但求值点必须一致。
\emph{核验：}矩阵式逐行就是两方程的链式求导。

\textbf{例4（Jacobian 退化）。}
\emph{信号：}\(F=u^2-x=0,G=v-y=0\)，原点。
\emph{分析与解答：}关于 \((u,v)\) 的 Jacobian 在原点为
\(\begin{vmatrix}0&0\\0&1\end{vmatrix}=0\)；\(u=\pm\sqrt{x}\) 多分支且在
\(x=0\) 导数发散，不能套隐函数组公式。
\emph{条件与易错：}退化只说明定理不适用，应进一步观察实际分支。
\emph{核验：}同一 \(x>0\) 有两个 \(u\)。

\textbf{失分防线、考场表达与标准。}
错误是把 Jacobian 的变量顺序换了却不追踪符号，或直接分别解函数造成漏支。
考场先写对目标变量的系数矩阵，再写右端并解线性系统。
基础过关会隐式求导；130 分以上会处理两个自变量；
满分能力能解释 Jacobian 非零为何保证局部唯一与可解。
\textbf{连接：}行列式结构连接线性方程组、空间曲线切向量与变量变换。
\end{knowledgeBox}

\studySubsection{calc-13-k147}{K147 齐次函数与 Euler 公式}
\knowledgeAnchor[MATH1-KN-CALC-0194]{齐次函数与Euler公式}

\begin{knowledgeBox}
\textbf{定位与前置桥。}
齐次性描述同时缩放所有自变量时函数按固定幂缩放；它通过
\knowledgeRef[MATH1-KN-CALC-0189]{链式法则}把缩放参数的导数传回各偏导：
\[
f(tx,ty)=t^m f(x,y).
\]
它是什么：尺度规律；为何需要：把抽象缩放关系转成偏导恒等式；怎样使用：
对 \(t\) 求导后令 \(t=1\)；最小例子 \(x^2+xy+y^2\) 为二次齐次函数。
若 \(f\) 可微且缩放过程留在定义域内，则 Euler 公式为
\[
xf_x+yf_y=mf.
\]

\textbf{方法选择：}先用同次缩放或各项总次数确认齐次性；已知缩放关系时
对参数 \(t\) 求导并令 \(t=1\)；已知偏导恒等式时再反向识别次数，
不能在非锥形定义域上无条件外推。

\textbf{例1（识别齐次次数）。}
\emph{信号：}\(f=(x^2+y^2)^{3/2}\)。
\emph{分析与解答：}
\[
f(tx,ty)=(t^2x^2+t^2y^2)^{3/2}=t^3f(x,y)
\]
（取 \(t>0\)），故次数为 \(3\)。
\emph{条件与易错：}实数幂下先限定缩放符号；不能只看括号外指数 \(3/2\)。
\emph{核验：}各基本量 \(x^2,y^2\) 同为二次。

\textbf{例2（Euler 公式化简）。}
\emph{信号：}\(f=x^3-2x^2y+xy^2\)。
\emph{分析与解答：}各项总次数均为 \(3\)，故不必展开偏导即可知
\[
xf_x+yf_y=3f.
\]
\emph{条件与易错：}每一项总次数必须一致。
\emph{核验：}直接计算
\(f_x=3x^2-4xy+y^2,\ f_y=-2x^2+2xy\) 后相加。

\textbf{例3（由缩放关系推导公式）。}
\emph{信号：}已知 \(f(tx,ty)=t^m f(x,y)\)。
\emph{分析与解答：}对 \(t\) 求导：
\[
xf_x(tx,ty)+yf_y(tx,ty)=mt^{m-1}f(x,y).
\]
令 \(t=1\) 即得 Euler 公式。
\emph{条件与易错：}左侧偏导在 \((tx,ty)\) 求值；不能先令 \(t=1\) 再求导。
\emph{核验：}代入二次多项式恢复 \(2f\)。

\textbf{例4（二阶扩展）。}
\emph{信号：}\(f\in C^2\) 为 \(m\) 次齐次函数。
\emph{分析与解答：}对 \(xf_x+yf_y=mf\) 分别关于 \(x,y\) 求导，再乘
\(x,y\) 相加，得到
\[
x^2f_{xx}+2xyf_{xy}+y^2f_{yy}=m(m-1)f.
\]
\emph{条件与易错：}交叉项出现两次；使用 \(f_{xy}=f_{yx}\) 需相应连续条件。
\emph{核验：}对 \(f=x^2+y^2,m=2\)，两边均为 \(2(x^2+y^2)\)。

\textbf{失分防线、考场表达与标准。}
错误是把次数不一致的和式判为齐次，或漏 Euler 公式右侧的 \(m\)。
考场先写缩放检验，再由 Euler 公式化简。
基础过关能判次数；130 分以上能从公式求抽象偏导组合；
满分能力能从缩放关系现场推导并得到二阶恒等式。
\textbf{连接：}Euler 公式本质上是沿径向缩放方向使用链式法则。
\end{knowledgeBox}

\studySubsection{calc-13-k148}{K148 方向导数定义}
\knowledgeAnchor[MATH1-KN-CALC-0195]{方向导数定义}

\begin{knowledgeBox}
\textbf{定位与前置桥。}
\knowledgeRef[MATH1-KN-CALC-0184]{偏导}只看坐标轴方向；
方向导数沿任意给定方向观察变化，而
\knowledgeRef[MATH1-KN-CALC-0188]{可微性定义检验}提醒我们逐方向存在仍不等于
统一线性近似。
单位向量是什么：长度为 \(1\) 的方向代表；为何需要：参数 \(t\) 才等于实际有向距离；
怎样使用：先单位化 \(\boldsymbol\ell\)，再研究一元函数
\(g(t)=f(P+t\boldsymbol\ell)\)；最小例子是
\(\boldsymbol a=(3,4)\) 对应单位方向 \((3/5,4/5)\)。
\[
D_{\boldsymbol\ell}f(P)
=\lim_{t\to0}\frac{f(P+t\boldsymbol\ell)-f(P)}t.
\]
方向导数存在不等于函数可微。

\textbf{方法选择：}给定非单位方向先单位化；已知函数可微时可转用 K149
的梯度点积公式；分段或非光滑点必须回到单变量差商定义，并比较正、负
参数及不同方向。

\textbf{例1（先单位化）。}
\emph{信号：}\(f=x^2+y^2\)，点 \((1,2)\)，方向 \((3,4)\)。
\emph{分析与解答：}\(\boldsymbol\ell=(3/5,4/5)\)。按可微公式或定义，
\[
D_{\boldsymbol\ell}f=(2,4)\cdot(3/5,4/5)=22/5.
\]
\emph{条件与易错：}直接点乘 \((3,4)\) 会多乘长度 \(5\)。
\emph{核验：}把 \((1,2)+t(3/5,4/5)\) 代入差商。

\textbf{例2（由方向角）。}
\emph{信号：}方向与 \(x\) 轴正向夹角为 \(\alpha\)。
\emph{分析与解答：}\(\boldsymbol\ell=(\cos\alpha,\sin\alpha)\)。对
\(f=xy\) 在 \((1,1)\)，
\[
D_{\boldsymbol\ell}f=\cos\alpha+\sin\alpha.
\]
\emph{条件与易错：}方向角自带单位向量，不需再次除长度。
\emph{核验：}\(\alpha=0,\pi/2\) 分别恢复 \(f_x,f_y\)。

\textbf{例3（特殊点按定义）。}
\emph{信号：}\(f(x,y)=\sqrt{x^2+y^2}\) 于原点。
\emph{分析与解答：}对任意单位向量 \(\boldsymbol\ell\)，
\[
\frac{f(t\boldsymbol\ell)-f(0)}t=\frac{|t|}{t},
\]
双侧极限不存在，故按双侧定义所有方向导数均不存在。
\emph{条件与易错：}若题目采用单侧方向导数约定会得到 \(1\)，必须服从题设口径；
本仓库默认双侧。
\emph{核验：}正负 \(t\) 给出相反值。

\textbf{例4（所有方向导数存在仍不可微）。}
\emph{信号：}
\[
f(x,y)=
\begin{cases}\dfrac{x^3}{x^2+y^2},&(x,y)\ne(0,0),\\0,&(0,0).\end{cases}
\]
\emph{分析与解答：}沿单位方向 \((a,b)\)，
\[
\frac{f(ta,tb)}t=\frac{a^3}{a^2+b^2}=a^3,
\]
故各方向导数存在；但方向值 \(a^3\) 不是关于 \((a,b)\) 的线性函数，
且 K140 已证该函数不可微。
\emph{条件与易错：}方向导数的逐方向存在没有给出统一线性近似。
\emph{核验：}若可微且梯度为 \((1,0)\)，方向导数应为 \(a\)，与 \(a^3\) 冲突。

\textbf{失分防线、考场表达与标准。}
错误是忘记单位化或把“方向导数存在”当“可微”。
考场先写单位方向，再写参数化差商；已知可微时才可直接用梯度点积。
基础过关会按定义求值；130 分以上能处理特殊点；
满分能力能用方向导数关于方向的非线性否定可微。
\textbf{连接：}可微时全部方向导数由
\knowledgeRef[MATH1-KN-CALC-0196]{梯度}统一编码。
\end{knowledgeBox}

\studySubsection{calc-13-k149}{K149 梯度及最大方向导数}
\knowledgeAnchor[MATH1-KN-CALC-0196]{梯度及最大方向导数}

\begin{knowledgeBox}
\textbf{前置短桥与核心。}
\knowledgeRef[MATH1-KN-CALC-0195]{方向导数}沿单位方向参数化；点积
\(\boldsymbol a\cdot\boldsymbol\ell
=\|\boldsymbol a\|\cos\theta\) 衡量向量在该方向的投影。
若 \(f\) 在点可微，则
\[
\nabla f=(f_x,f_y),\qquad
D_{\boldsymbol\ell}f=\nabla f\cdot\boldsymbol\ell.
\]
由 Cauchy--Schwarz 不等式，最大方向导数为 \(\|\nabla f\|\)，方向为非零梯度的
单位方向；最小值为其负值。梯度为零只说明一阶变化率全为零，不自动说明极值。

\textbf{方法选择：}指定方向就计算梯度与单位向量的点积；求最大、最小
方向导数用 Cauchy--Schwarz；若梯度为零，只能得到所有一阶方向导数为零，
极值还需另用高阶或定义判定。

\textbf{例1（指定方向与最大方向）。}
\emph{信号：}\(f=x^2+xy+y^2\)，点 \((1,0)\)。
\emph{分析与解答：}\(\nabla f=(2x+y,x+2y)\)，故该点为 \((2,1)\)。
沿 \((1,1)\) 的单位方向有 \(D=3/\sqrt2\)；最大方向导数为 \(\sqrt5\)，
方向为 \((2/\sqrt5,1/\sqrt5)\)。
\emph{条件与易错：}指定方向先单位化。
\emph{核验：}\(3/\sqrt2<\sqrt5\) 符合最大值界。

\textbf{例2（最速下降）。}
\emph{信号：}温度场 \(T=x^2+4y^2\)，点 \((1,1)\)。
\emph{分析与解答：}\(\nabla T=(2,8)\)，最速下降方向为
\((-1/\sqrt{17},-4/\sqrt{17})\)，最小方向导数为 \(-2\sqrt{17}\)。
\emph{条件与易错：}问“下降最快”要取负梯度，不是梯度。
\emph{核验：}方向向量长度为 \(1\)。

\textbf{例3（梯度为零不判极值）。}
\emph{信号：}\(f=x^2-y^2\) 于原点。
\emph{分析与解答：}\(\nabla f(0,0)=0\)，所以由可微公式各方向一阶导数均为零；
但沿 \(y=0\) 函数为正，沿 \(x=0\) 为负，原点是鞍点而非极值。
\emph{条件与易错：}梯度为零只是可微内点极值的必要条件。
\emph{核验：}两条路径直接给出异号增量。

\textbf{例4（等值线法向）。}
\emph{信号：}等值线 \(x^2+xy+y^2=3\) 在 \((1,1)\)。
\emph{分析与解答：}令 \(F=x^2+xy+y^2\)，
\(\nabla F(1,1)=(3,3)\)，它垂直该点等值线；切线方向可取 \((1,-1)\)。
\emph{条件与易错：}需确认梯度非零且点在等值线上。
\emph{核验：}\((3,3)\cdot(1,-1)=0\)。

\textbf{失分防线、考场表达与标准。}
错误是未单位化、把最大方向导数写成梯度本身，或由零梯度直接判极值。
考场先求并代点得到梯度，再用点积和模。
基础过关会算指定方向；130 分以上能给最速方向与数值；
满分能力能用向量几何解释等值线法向和零梯度退化。
\textbf{连接：}三元函数梯度是隐式曲面的法向量。
\end{knowledgeBox}

\studySubsection{calc-13-k150}{K150 曲面切平面与法线}
\knowledgeAnchor[MATH1-KN-CALC-0197]{曲面切平面与法线}

\begin{knowledgeBox}
\textbf{定位与前置桥。}
平面方程可由一个点和法向量确定；梯度垂直函数的等值面。
它是什么：\(\nabla F(P)\) 是隐式曲面 \(F=0\) 在正则点 \(P\) 的法向量；
为何需要：直接写切平面和法线；怎样使用：先验点在曲面上且梯度非零；
最小模型是平面 \(F=ax+by+cz-d\)，其梯度正是固定法向量。
\[
\text{切平面： }\nabla F(P)\cdot(\boldsymbol r-P)=0,\qquad
\text{法线： }\boldsymbol r=P+t\nabla F(P).
\]

\textbf{方法选择：}隐式曲面直接取 \(\nabla F\) 作法向量；显式曲面
\(z=f(x,y)\) 先改写为 \(F=f-z\)。落笔前先验点在曲面上且
\(\nabla F(P)\ne\boldsymbol0\)，再分别写点法式切平面与参数式法线。

\textbf{例1（隐式球面）。}
\emph{信号：}\(x^2+y^2+z^2=3\)，点 \(P=(1,1,1)\)。
\emph{分析与解答：}\(\nabla F(P)=(2,2,2)\)，故切平面
\[
(x-1)+(y-1)+(z-1)=0
\]
即 \(x+y+z=3\)；法线为
\((x,y,z)=(1,1,1)+t(1,1,1)\)。
\emph{条件与易错：}先核验 \(P\) 在球面；法向量可约去非零公因子。
\emph{核验：}球心到 \(P\) 的半径方向与法线一致。

\textbf{例2（显式曲面）。}
\emph{信号：}\(z=x^2+y^2\)，点 \((1,1,2)\)。
\emph{分析与解答：}令 \(F=x^2+y^2-z\)，
\(\nabla F=(2x,2y,-1)\)，在点为 \((2,2,-1)\)。切平面为
\[
2(x-1)+2(y-1)-(z-2)=0.
\]
\emph{条件与易错：}显式曲面必须改写为 \(F=0\)，梯度第三分量是 \(-1\)。
\emph{核验：}等价写成 \(z-2=2(x-1)+2(y-1)\)，与全微分一致。

\textbf{例3（切平面平行条件）。}
\emph{信号：}在曲面 \(z=x^2+y^2\) 上找切平面平行于
\(2x+2y-z=0\) 的点。
\emph{分析与解答：}曲面法向量 \((2x,2y,-1)\) 应与 \((2,2,-1)\) 平行。
第三分量迫使比例系数为 \(1\)，故 \(x=y=1,z=2\)。
\emph{条件与易错：}最后必须用曲面方程补出 \(z\) 并核验。
\emph{核验：}所得点正是例2的点。

\textbf{例4（正则条件失效）。}
\emph{信号：}锥面 \(F=x^2+y^2-z^2=0\) 于原点。
\emph{分析与解答：}\(\nabla F(0,0,0)=0\)，梯度法无法给出唯一法向量；
锥顶确实没有唯一切平面。
\emph{条件与易错：}把零向量代入平面公式会得到 \(0=0\)，不是切平面。
\emph{核验：}沿不同母线靠近锥顶，候选切平面方向不同。

\textbf{失分防线、考场表达与标准。}
错误是漏代点、把切向量当法向量，或在梯度为零时仍套公式。
考场顺序固定为“验点—求梯度并代点—验非零—写平面/法线”。
基础过关会处理显式与隐式曲面；130 分以上能解平行、垂直条件；
满分能力能识别正则点假设并用全微分、几何两路核验。
\textbf{连接：}本节把
\knowledgeRef[MATH1-KN-CALC-0192]{隐函数}与
\knowledgeRef[MATH1-KN-CALC-0196]{梯度}汇合为几何应用。
\end{knowledgeBox}
